\documentclass[11pt,letterpaper]{article}
\usepackage[margin=2.5cm]{geometry}
\usepackage[utf8]{inputenc}
\usepackage[T1]{fontenc}
\usepackage{lmodern}
\usepackage{amsmath,amssymb,amsthm,mathtools,bm}
\usepackage{enumitem}
\usepackage{booktabs}
\usepackage{tikz}
\usetikzlibrary{arrows.meta,patterns,positioning,calc}
\usepackage{pgfplots}
\pgfplotsset{compat=1.16}
\usepackage{titlesec}
\usepackage[hidelinks]{hyperref}

\setlength{\parindent}{0pt}
\setlength{\parskip}{0.5em}
\titleformat{\section}{\large\bfseries}{\thesection.}{0.6em}{}
\titleformat{\subsection}{\normalsize\bfseries}{\thesubsection.}{0.6em}{}

\theoremstyle{definition}
\newtheorem{definition}{Definition}[section]
\newtheorem{example}{Example}[section]
\newtheorem{exercise}{Exercise}[section]
\newtheorem{theorem}{Theorem}[section]
\newtheorem{remark}{Remark}[section]

\newcommand{\R}{\mathbb{R}}
\newcommand{\N}{\mathbb{N}}
\newcommand{\Z}{\mathbb{Z}}
\newcommand{\U}{\mathcal{U}}
\newcommand{\E}{\mathbb{E}}
\newcommand{\Var}{\operatorname{Var}}
\newcommand{\Cov}{\operatorname{Cov}}
\newcommand{\Normal}{\mathcal{N}}
\newcommand{\dd}{\,\mathrm{d}}
\newcommand{\T}{^{\top}}
\renewcommand{\P}{P}
\newcommand{\set}[1]{\left\{#1\right\}}
\newcommand{\given}{\mid}
\newcommand{\note}[1]{\textit{#1}}

\title{Lecture 4: Classical Inference}
\author{}\date{}

\begin{document}
\maketitle

\section*{Non-Bayesian methods of doing statistics}
Here, I am going to talk about some of the most used methods for doing non-Bayesian statistics.

%=====================================================
\section{Maximum likelihood estimation (MLE)}
It is a form of point estimation. The goal is to find the parameter $\theta$:
\begin{align*}
\hat\theta_{\text{mle}} &= \arg\max_{\theta}\, p(Y \mid X, \theta) && \text{(estimator, } \hat\theta(Y)\text{)},\\
\hat\theta_{\text{mle}} &= \arg\max_{\theta}\, p(\vec y \mid \vec x, \vec\theta) && \text{(estimate)},
\end{align*}
where
\[
p(\vec y \mid \vec x, \vec\theta) = p(y_1, \dots, y_N \mid x_1, \dots, x_N, \theta_1, \dots, \theta_K).
\]
The maximum likelihood estimate is the value of $\theta$ that makes the data $\vec y$ most likely.

When the samples are independent and identically distributed (iid), then
\[
p(\vec y \mid \vec x, \vec\theta) = p(y_1 \mid x_1, \vec\theta)\,p(y_2 \mid x_2, \vec\theta)\cdots p(y_N \mid x_N, \vec\theta) = \prod_{i=1}^{N} p(y_i \mid x_i, \vec\theta).
\]
Normally, we use the log-likelihood:
\[
\log p(\vec y \mid \vec x, \vec\theta) = \sum_{i=1}^{N}\log p(y_i \mid x_i, \vec\theta),
\]
\[
\hat\theta = \arg\max_{\theta}\sum_{i=1}^{N}\log p(y_i \mid x_i, \vec\theta)
\qquad\text{or}\qquad
\hat\theta = \arg\min_{\theta}\left[-\sum_{i=1}^{N}\log p(y_i \mid x_i, \vec\theta)\right].
\]
Alternatively, one may want to maximize the joint likelihood:
\[
\hat\theta = \arg\min_{\theta}\left[-\sum_{i=1}^{N}\log p(y_i, x_i \mid \vec\theta)\right].
\]

\begin{example}[Bernoulli distribution]
\[
p(Y = y_i \mid \theta) = \theta^{y_i}(1-\theta)^{1-y_i}, \qquad y_i \in \set{0,1}.
\]
In that case, we aim to find
\[
\hat\theta = \arg\max_{\theta}\sum_{i=1}^{N}\log p(Y = y_i \mid \theta),
\]
\[
L(\theta) = \sum_{i=1}^{N}\log\left[\theta^{y_i}(1-\theta)^{1-y_i}\right] = \sum_{i=1}^{N}\left[\new{y_i\log\theta + (1-y_i)\log(1-\theta)}\right].
\]
\[
\frac{\partial L}{\partial\theta} = \sum_{i=1}^{N}\left(y_i\new{\frac{1}{\theta}} - (1-y_i)\new{\frac{1}{1-\theta}}\right)
= \sum_{i=1}^{N}\left(\frac{y_i}{\theta} \new{{}+ \frac{y_i - 1}{1-\theta}}\right)
= \sum_{i=1}^{N}\frac{\new{y_i(1-\theta) + \theta(y_i - 1)}}{\new{\theta(1-\theta)}} = 0.
\]
Thus
\[
\sum_{i}\left(\new{y_i - y_i\theta + \theta y_i - \theta}\right) = 0
\ \Rightarrow\
\sum_{i=1}^{N} y_i - \theta\underbrace{\sum_{i=1}^{N} 1}_{\new{N}} = 0
\ \Rightarrow\
\boxed{\hat\theta = \frac{1}{N}\sum_{i=1}^{N} y_i}\quad\text{(MLE)}.
\]
\end{example}

%=====================================================
\section{Confidence intervals}
A $100(1-\alpha)\%$ \emph{confidence interval} for a parameter $\theta$ is any interval $I = (l, u)$ such that
\[
p(\theta \in I) = 1 - \alpha.
\]
Here $\theta$ is fixed and the interval $I$ is random (it is computed from the sample), so the probability refers to $I$.
To study confidence intervals (CI) we need sampling distributions.

\begin{theorem}
Let a population be normally distributed. This is, there is a random variable $X_i \sim \Normal(\mu, \sigma^2)$. Then, if a sample of size $n$ is drawn from the population, the sampling mean $\bar X = \frac1n\sum_i X_i$ satisfies
\[
\bar X \sim \Normal\left(\mu, \frac{\sigma^2}{n}\right),
\qquad
Z = \frac{\bar X - \mu}{\sigma/\sqrt n} \sim \Normal(0,1),
\]
with $\sigma/\sqrt n$ the \emph{standard error}.
\end{theorem}

So, we aim to set a CI to estimate something about a population parameter:
\[
p\big(|\bar X - \mu| < x\big) = 1 - \alpha
\ \Rightarrow\
p\left(\frac{|\bar X - \mu|}{\new{\sigma/\sqrt n}} < \frac{x}{\new{\sigma/\sqrt n}}\right) = 1 - \alpha,
\]
\[
p\big(|Z| < z_\alpha\big) = 1 - \alpha
\ \Rightarrow\
p(\new{-z_\alpha < Z < z_\alpha}) = 1 - \alpha.
\]
Then
\[
-z_\alpha < \frac{\bar X - \mu}{\sigma/\sqrt n} < z_\alpha
\ \Rightarrow\
-z_\alpha\new{\frac{\sigma}{\sqrt n}} < \bar X - \mu < z_\alpha\new{\frac{\sigma}{\sqrt n}}.
\]
Multiplying by minus,
\[
\new{z_\alpha}\frac{\sigma}{\sqrt n} \mathrel{\new{>}} \new{\mu - \bar X} \mathrel{\new{>}} \new{-z_\alpha}\frac{\sigma}{\sqrt n}
\ \Rightarrow\
\new{\bar X} - z_\alpha\frac{\sigma}{\sqrt n} < \mu < \new{\bar X} + z_\alpha\frac{\sigma}{\sqrt n}.
\]
The interval is
\[
\left[\bar X - z_\alpha\frac{\sigma}{\sqrt n},\ \bar X + z_\alpha\frac{\sigma}{\sqrt n}\right].
\]
We only need to find $z_\alpha$:
\[
p(-z_\alpha < Z < z_\alpha) = 1 - \alpha
\ \Rightarrow\
\new{p(Z \le z_\alpha) - p(Z \le -z_\alpha)} = 1 - \alpha,
\]
where $p(Z \le -z_\alpha) = p(Z > z_\alpha) = 1 - p(Z \le z_\alpha)$. Hence
\[
\new{2}p(Z \le z_\alpha) \new{{}- 1} = 1 - \alpha
\ \Rightarrow\
p(Z \le z_\alpha) = \new{1 - \frac{\alpha}{2}}.
\]
(This value is often written $z_{\alpha/2}$; e.g.\ $z_\alpha = 1.96$ for $\alpha = 0.05$.)

\begin{example}
See Pg.~36, \emph{Hypothesis testing 2}.
\end{example}

\begin{theorem}
Let $X_i$ be random variables, $X_i \sim \Normal(\mu, \sigma^2)$, and $S^2$ be the sampling variance. Then
\[
T = \frac{\bar X - \mu}{S/\sqrt n} \sim t(n-1).
\]
\end{theorem}
In the same way,
\[
p(-t_\alpha < T < t_\alpha) = 1 - \alpha
\ \Rightarrow\
\left[\bar X - t_\alpha\frac{S}{\sqrt n},\ \bar X + t_\alpha\frac{S}{\sqrt n}\right].
\]

%=====================================================
\section{Hypothesis testing}
A \emph{hypothesis} is a statement about population parameters. We have two hypotheses:
\begin{itemize}[nosep]
\item null hypothesis $H_0$,
\item alternative hypothesis $H_1$.
\end{itemize}
$H_0$ is assumed to be true, and it is either rejected or not rejected. Null hypothesis: no change.

$\alpha$: significance level.
\begin{align*}
p(\text{reject } H_0 \mid H_0 \text{ true}) &= \alpha,\\
p(\text{not reject } H_0 \mid H_0 \text{ false}) &= \beta,\\
p(\text{reject } H_0 \mid H_0 \text{ false}) &= 1 - \beta.
\end{align*}

We need a statistic and a rejection region. We can have
\[
\begin{cases}
H_0: \theta = \theta_0\\
H_1: \theta > \theta_0
\end{cases}
\qquad
\begin{cases}
H_0: \theta = \theta_0\\
H_1: \theta < \theta_0
\end{cases}
\qquad
\begin{cases}
H_0: \theta = \theta_0\\
H_1: \theta \neq \theta_0
\end{cases}
\]
$H_1$ is contradictory to $H_0$.

How to find the rejection region? We set
\[
p(\text{type I error}) = \alpha.
\]

\begin{example}[Cigarettes, Pg.~51]
$n = 30$, $\alpha = 0.05$, $X \sim \Normal(\mu, \sigma^2)$,
\[
H_0: \mu \le 1.5\ \text{mg}, \qquad H_1: \mu > 1.5\ \text{mg}.
\]
I need a statistic to compare to $\mu$: that is $\bar X$.
\begin{align*}
p(\bar X > x_c) &= 0.05,\\
p\left(\frac{\bar X - \mu}{\sigma/\sqrt n} > z\right) &= 0.05,\\
1 - p\left(\frac{\bar X - \mu}{\sigma/\sqrt n} \le z\right) &= 0.05,\\
p\left(\frac{\bar X - \mu}{\sigma/\sqrt n} \le z\right) &= 0.95 \ \Rightarrow\ z = 1.64.
\end{align*}
Rejection region:
\[
\boxed{\bar X > \mu + 1.64\frac{\sigma}{\sqrt n}}
\]
with $\mu = 1.5$ mg.
\end{example}

\end{document}
